\documentclass[12pt]{article}
\usepackage[T1]{fontenc}
\usepackage[utf8]{inputenc}
\usepackage{lmodern}
\usepackage{textcomp}
\usepackage{enumitem}
\usepackage{geometry}
\geometry{margin=1in}
\begin{document}
\begin{center}
Answer Key - Computer Science - Graduate Level Assessment 16 \
\small (Answers are ordered exactly as in the question paper.)
\end{center}

\section*{Multiple Choice}
\begin{enumerate}[label=\arabic*.]
\item \textbf{Answer:} A
\\ \textit{Options:} A) WEP; B) WPA; C) WPA 2; D) WPA 3
\\ \textit{Note:} Correct option: A - WEP
\item \textbf{Answer:} C
\\ \textit{Options:} A) Love Bug; B) SQL Slammer; C) Morris Worm; D) Code Red
\\ \textit{Note:} Correct option: C - Morris Worm
\item \textbf{Answer:} C
\\ \textit{Options:} A) boundary hacks; B) memory checks; C) boundary checks; D) buffer checks
\\ \textit{Note:} Correct option: C - boundary checks
\item \textbf{Answer:} D
\\ \textit{Options:} A) DNSlookup; B) Whois; C) Nslookup; D) IP Network Browser
\\ \textit{Note:} Correct option: D - IP Network Browser
\item \textbf{Answer:} A
\\ \textit{Options:} A) one in which each main memory word can be stored at any of 3 cache locations; B) effective only if 3 or fewer processes are running alternately on the processor; C) possible only with write-back; D) faster to access than a direct-mapped cache
\\ \textit{Note:} Correct option: A - one in which each main memory word can be stored at any of 3 cache locations
\end{enumerate}

\section*{True / False}
\begin{enumerate}[label=\arabic*.]
\setcounter{enumi}{5}
\item \textbf{Answer:} False
\\ \textit{Options:} A) True; B) False
\\ \textit{Note:} LLM-generated false statement. Correct answer: 'No, an attacker can simply guess the tag for messages'.
\item \textbf{Answer:} True
\\ \textit{Options:} A) True; B) False
\\ \textit{Note:} LLM-generated true statement based on MCQ.
\item \textbf{Answer:} False
\\ \textit{Options:} A) True; B) False
\\ \textit{Note:} LLM-generated false statement. Correct answer: 'Assume that the query is not satisfiable and stop executing the path'.
\item \textbf{Answer:} True
\\ \textit{Options:} A) True; B) False
\\ \textit{Note:} LLM-generated true statement based on MCQ.
\item \textbf{Answer:} False
\\ \textit{Options:} A) True; B) False
\\ \textit{Note:} LLM-generated false statement. Correct answer: 'Key exchange'.
\end{enumerate}

\section*{Long Form Response}
\begin{enumerate}[label=\arabic*.]
\setcounter{enumi}{10}
\item \textbf{Answer:} def sum\_positivenum(nums): | return sum(sum\_positivenum)
\\ \textit{Note:} Students should provide a detailed explanation referencing algorithm steps.
\item \textbf{Answer:} def check\_subset\_list(list 1, list 2): | return exist
\\ \textit{Note:} Students should provide a detailed explanation referencing algorithm steps.
\end{enumerate}

\vfill
\noindent\textit{End of Answer Key}
\end{document}